\section{Leadership \& Team Reflection}\label{sec:team-reflection}

\subsection{Leadership Structure and Roles}\label{sec:roles}
Given the purely analytical nature of this project and our large team size of six members, we determined that a highly compartmentalized structure would be inefficient. Instead, we adopted a \textbf{Paired-Leadership Model}. We assigned two members to each of the three mandatory roles, allowing for deeper collaborative analysis and brainstorming without relying on technical implementation tasks:
\begin{itemize}
    \item \textbf{Co-Project Leaders (2 Members):} Responsible for defining the overall project scope, steering the brainstorming sessions (\hyperref[sec:brainstorming]{Stage~1}), and structuring the innovative solution (\hyperref[sec:solution-design]{Stage~3}). They ensured the team's proposed ``Fair Trade AI'' framework remained realistic and cohesive.
    \item \textbf{Co-Ethics \& Legal Officers (2 Members):} Tasked with conducting the critical dilemma analysis (\hyperref[sec:critical-analysis]{Stage~2}) and researching professional codes of conduct. They acted as the moral compass of the group, focusing entirely on the societal impacts, copyright laws, and the ethical justification of our solution (\hyperref[sec:ethics-legal]{Stage~4}).
    \item \textbf{Co-Communication Leads (2 Members):} Managed the documentation and presentation aspects. They were responsible for synthesizing the analytical work of the entire team into a single, cohesive LaTeX report, ensuring a consistent professional tone, and preparing the structure for the final oral presentation.
\end{itemize}

\subsection{Conflict Resolution and Decision-Making Process}\label{sec:conflict-resolution}
Because our project focused on conceptual analysis rather than writing code, our conflicts were primarily ideological. During the dilemma analysis phase, our team experienced a significant debate regarding how strict our proposed solution should be.

One pair argued that to truly protect intellectual property, Generative AI models should be temporarily banned from commercial use until global copyright laws catch up. The other pair argued that a blanket ban is not a realistic engineering solution and that we must design a path forward (the Opt-in and Watermarking system).

To resolve this, we used a consensus-building approach, evaluating both arguments against the course's definition of professional engineering practice. We decided that engineers must solve problems proactively rather than just halting progress. We agreed on the ``Fair Trade AI'' model, combining technological innovation with strict ethical boundaries, thereby resolving the conflict through critical thinking rather than compromise.

\subsection{Reflection on Team Dynamics and Professional Growth}\label{sec:team-dynamics}
\begin{itemize}
    \item \textbf{Team Strengths:} The paired-leadership structure was our greatest asset. By having two minds focused on the ethical analysis and two on the solution design, we were able to brainstorm a much wider range of scenarios and stakeholder impacts than a single person could have achieved.
    \item \textbf{Areas for Improvement:} Initially, having six members inputting ideas led to a lack of focus. Our early drafts read like a disjointed collection of opinions rather than a unified professional report. We learned that intensive analytical work requires strict synthesis, forcing us to conduct multiple group review sessions to align our arguments.
    \item \textbf{Distribution of Contributions:} The analytical workload was distributed equitably across the pairs:
    \begin{itemize}
        \item \todo{These pairs do not match the Contribution column of the member table. Make them agree.}\textbf{Members 1 \& 2 (Project Management):} Led the SCAMPER/Mind-mapping sessions, formulated the core problem framing, and structured the proposed Fair Trade AI solution and its trade-offs.
        \item \textbf{Members 3 \& 4 (Ethical Analysis):} Conducted the deep dive into IEEE/ACM guidelines, authored the dilemma analysis regarding deepfakes vs.\ innovation, and explicitly defined the absolute boundaries of engineering responsibility.
        \item \textbf{Members 5 \& 6 (Documentation \& Synthesis):} Managed the LaTeX source code, integrated the diverse written sections into a unified narrative, and finalized the report formatting to meet the strict 10--15 page requirement.
    \end{itemize}
\end{itemize}
